\documentclass[11pt,a4paper]{ctexart}
\input{/Users/shaoyikai/Desktop/考研/22408-notes-skill/assets/note-preamble.tex}
\renewcommand{\NoteFooter}{22408 一页纸笔记 · 数学二}
\linespread{1.14}
\begin{document}
\notetitle{反常积分敛散性判别}{数学二}{高等数学·定积分}
\sechead{一、定义与判敛总纲}

\begin{itemize}

\item \textbf{无穷区间}：$\displaystyle\int_a^{+\infty}f\,\mathrm{d}x=\lim_{b\to+\infty}\int_a^{b}f\,\mathrm{d}x$；\textbf{瑕积分}（$x=c$ 处 $f$ 无界）：$\displaystyle\int_a^{c}f\,\mathrm{d}x=\lim_{t\to c^{-}}\int_a^{t}f\,\mathrm{d}x$。

\item \textbf{收敛}：该极限存在且有限。$\displaystyle\int_{-\infty}^{+\infty}f\,\mathrm{d}x$ 要求两段\textbf{同时收敛}，不能用对称抵消（如 $\displaystyle\int_{-\infty}^{+\infty}x\,\mathrm{d}x$ 发散）。

\end{itemize}

\notebox{判敛只看"速度"：$x\to+\infty$ 时 $f\to0$ 的快慢；$x\to c$ 时 $f\to\infty$ 的快慢。}

\sechead{二、核心标尺：$p$ 积分（必背）}

\begin{center}\small
\begin{tabularx}{\linewidth}{@{}>{\raggedright\arraybackslash}X>{\raggedright\arraybackslash}X>{\raggedright\arraybackslash}X@{}}
\toprule
类型 & 基准积分 & 收敛条件 \\
\midrule
无穷区间 & $\displaystyle\int_1^{+\infty}\frac{\mathrm{d}x}{x^{p}}$ & $p>1$ 收敛；$p\le1$ 发散 \\
瑕积分（瑕点 $a$） & $\displaystyle\int_a^{b}\frac{\mathrm{d}x}{(x-a)^{p}}$ & $p<1$ 收敛；$p\ge1$ 发散 \\
\bottomrule
\end{tabularx}
\end{center}

\warnline{两个结论的 $p$ 方向\textbf{相反}：无穷远处要衰减得够快（$p>1$），瑕点处要爆得不够快（$p<1$）。}

\sechead{三、比较判别法}

\textbf{无穷区间}（$f,g\ge0$，$x\ge a$）：取 $\displaystyle\lim_{x\to+\infty}\frac{f(x)}{g(x)}=A$，则

\begin{itemize}

\item $0<A<+\infty$：$f$ 与 $g$ \textbf{同敛散}；

\item $A=0$：$g$ 收敛 $\Rightarrow$ $f$ 收敛（大收则小收）；

\item $A=+\infty$：$g$ 发散 $\Rightarrow$ $f$ 发散。

\end{itemize}

\textbf{瑕积分}（$x\to a^{+}$，$f,g\ge0$）：取 $\displaystyle\lim_{x\to a^{+}}\frac{f(x)}{g(x)}=A$，结论与上完全相同。

\warnline{比较法的标准动作是\textbf{抓主部}：$x\to+\infty$ 时只留增长最快的项，如 $\dfrac{x+1}{x^{3}+2}\sim\dfrac{1}{x^{2}}$；$x\to a$ 时只留最低阶项，如 $x\to1$ 时 $\ln x\sim x-1$。}

\sechead{四、收敛的必要条件与典型例题（本题即例 1）}

\notebox{必要条件：若 $\displaystyle\int_a^{+\infty}f\,\mathrm{d}x$ 收敛，则 $\displaystyle\lim_{x\to+\infty}f(x)=0$。反之不成立（$\displaystyle\int_1^{+\infty}\frac{\mathrm{d}x}{x}$ 发散但 $\frac1x\to0$），故它只能\textbf{快速判发散}。}

\textbf{例 1} $\displaystyle\int_2^{+\infty}\frac{1}{\sqrt{x}}\,\mathrm{d}x$：无穷区间 $p$ 积分，$p=\frac12\le1$，\textbf{发散}。直接算可验证：$\displaystyle\int_2^{b}x^{-1/2}\mathrm{d}x=2\sqrt{b}-2\sqrt2\to+\infty$。

\textbf{例 2} $\displaystyle\int_2^{+\infty}\frac{\ln x}{x}\,\mathrm{d}x$ 与 $\displaystyle\int_2^{+\infty}\frac{\mathrm{d}x}{x\ln x}$：均令 $t=\ln x$，分别化为 $\displaystyle\int_{\ln2}^{+\infty}t\,\mathrm{d}t$ 与 $\displaystyle\int_{\ln2}^{+\infty}\frac{\mathrm{d}t}{t}$，都\textbf{发散}。

\textbf{例 3} $\displaystyle\int_2^{+\infty}\frac{x}{e^{x}}\,\mathrm{d}x$：因 $\displaystyle\lim_{x\to+\infty}x^{2}\cdot\frac{x}{e^{x}}=0$，而 $\displaystyle\int_2^{+\infty}\frac{\mathrm{d}x}{x^{2}}$ 收敛，由比较法（$A=0$）知原积分\textbf{收敛}。

\notebox{快速判据：指数级衰减（$e^{-x}$、$e^{-x^{2}}$）必收敛；对数因子 $\ln x$ 增长最慢，一般不改变敛散性，但会拖慢幂函数，需换元或比较验证。}

\sechead{五、易错点}

\warnline{① 无穷区间判据是 $p>1$ 收敛，勿与瑕积分的 $p<1$ 记混，这是最高频失分点。}

\warnline{② 比较法必须先确认被积函数\textbf{非负}（或对其取绝对值），有正有负时不能直接比。}

\warnline{③ 瑕点要主动找：分母零点、$\ln x$ 在 $x=1$ 处、$\tan x$ 的间断点都可能是瑕点；区间内有瑕点必须拆开逐段判定。}

\warnline{④ $\dfrac{1}{x\ln x}$、$\dfrac{1}{x\ln^{2}x}$ 不能当 $p$ 积分直接套，须换元 $t=\ln x$ 化成 $p$ 积分；$\displaystyle\int_1^{+\infty}\frac{\ln x}{x^{p}}\,\mathrm{d}x$ 仍按 $p>1$ 判收敛，对数因子不会把门槛抬到 $p>2$。}

\end{document}
